%% ==================== 附录固定说明 模板 (v1.5.6, BZD 数模社 2026 规范) ====================
%%
%% 借鉴 BZD 数模社 附录模板: A 支撑材料清单 + B 软件环境 + C/D/E/F 源程序 + 补充
%%
%% 2026 规范 (第五/十一条 + 2026-08-26 官方答疑 Q3):
%% - 附录首页应是支撑材料清单 (file list)
%% - 附录 1: 程序使用说明
%% - 附录 2: 支撑材料说明
%% - 附录 3: 核心源代码 (推荐)
%% - 漏写 = 0 分 (评审硬性要求)
%%
%% 跑题用户必须:
%% 1. 删掉所有 % 注释
%% 2. 替换【】占位符为自己的实际文件清单和软件版本
%% 3. 删 下面的 \AppendixSlot{...} 行 (check 15 抓这个 inline 【】)

\section*{附录固定说明}

% ====== 附录 A: 支撑材料文件列表 (必填) ======
\subsection*{附录 A \quad 支撑材料文件列表}

\begin{table}[H]
  \centering
  \caption{支撑材料文件列表}
  \label{tab:appendix_files}
  \begin{tabular}{clll}
    \toprule
    \textbf{序号} & \textbf{文件名称} & \textbf{格式} & \textbf{主要内容} \\
    \midrule
    1 & \texttt{求解/问题1/问题1\_xxx.py} & PY & 问题一数据处理与模型求解 \\
    2 & \texttt{求解/问题1/图片/*.png} & PNG & 问题一输出图表 \\
    3 & \texttt{求解/问题1/结果/*.csv} & CSV & 问题一中间结果 \\
    4 & \texttt{求解/问题2/问题2\_yyy.py} & PY & 问题二预测与时点优化 \\
    5 & \texttt{求解/问题3/问题3\_zzz.py} & PY & 问题三分层优化 \\
    6 & \texttt{求解/问题4/问题4\_ddd.py} & PY & 问题四集成决策 \\
    7 & \texttt{求解/共享/constants.py} & PY & 跨问题物理常量 \\
    8 & \texttt{数据/原始数据.xlsx} & XLSX & 题目附件 (如有补充数据) \\
    9 & \texttt{AI工具使用详情.pdf} & PDF & AI 工具使用情况详细说明 \\
    \bottomrule
  \end{tabular}
\end{table}

% ====== 附录 B: 软件环境与求解工具 (必填) ======
\subsection*{附录 B \quad 软件环境与求解工具}

本文使用 \textbf{Python 3.12} 完成数据处理与模型求解, 主要程序库及版本:
\begin{itemize}
  \item numpy 1.26 (数值计算)
  \item pandas 2.0 (数据处理)
  \item scipy 1.11 (科学计算, 含 optimize 优化)
  \item scikit-learn 1.4 (机器学习, 含线性回归/随机森林/Isolation Forest)
  \item matplotlib 3.8 (可视化)
  \item PuLP 2.7 (线性规划, 如有优化问题)
\end{itemize}

使用 \textbf{Overleaf + xelatex (TeX Live 2024)} 完成论文排版, 使用 \textbf{TikZ} 绘制整体求解框架图。

% ====== 附录 C-F: 各问题源程序说明模板 (简化) ======
\subsection*{附录 C \quad 问题一完整源程序}

\begin{itemize}
  \item \textbf{程序名称}: 问题1\_xxx.py
  \item \textbf{编程语言及版本}: Python 3.12
  \item \textbf{对应问题}: 问题一
  \item \textbf{程序作用}: 【数据处理 / 模型求解 / 检验 / 绘图】
  \item \textbf{输入文件}: 【原始数据.xlsx 的 Sheet1】
  \item \textbf{输出文件}: 【问题1\_图片/*.png, 问题1\_结果/*.csv】
  \item \textbf{依赖环境}: numpy, pandas, scikit-learn
  \item \textbf{运行顺序}: \texttt{python 问题1/问题1\_xxx.py}
  \item \textbf{正文对应位置}: 第五章 5.1 节
\end{itemize}

完整源程序:
\begin{verbatim}
# 完整可运行的 Python 代码 (贴核心 30-50 行)
import numpy as np
import pandas as pd
from sklearn.ensemble import RandomForestRegressor

# 1. 加载数据
data = pd.read_excel('../../数据/原始数据.xlsx', sheet_name=0)

# 2. 数据预处理
X = data[['特征1', '特征2', '特征3']].values
y = data['目标变量'].values

# 3. 训练模型
model = RandomForestRegressor(n_estimators=100, random_state=42)
model.fit(X, y)

# 4. 预测与评估
y_pred = model.predict(X)
r2 = model.score(X, y)
print(f'R² = {r2:.4f}')
\end{verbatim}

\subsection*{附录 D \quad 问题二完整源程序}

\textbf{程序名称}: 问题2\_yyy.py

\textbf{编程语言及版本}: Python 3.12

\textbf{对应问题}: 问题二

\textbf{程序作用}: 【数据处理 / 模型求解 / 检验 / 绘图】

\textbf{输入文件}: 【原始数据.xlsx 的 Sheet1】+ 【问题一输出结果】

\textbf{输出文件}: 【问题2\_图片/*.png, 问题2\_结果/*.csv】

\textbf{依赖环境}: numpy, pandas, scikit-learn, scipy

\textbf{运行顺序}: \texttt{python 问题2/问题2\_yyy.py}

\textbf{正文对应位置}: 第五章 5.2 节

% 完整可运行代码见支撑材料对应文件

\subsection*{附录 E \quad 问题三完整源程序}

\textbf{程序名称}: 问题3\_zzz.py

\textbf{编程语言及版本}: Python 3.12

\textbf{对应问题}: 问题三

\textbf{程序作用}: 【数据处理 / 模型求解 / 检验 / 绘图】

\textbf{输入文件}: 【原始数据.xlsx 的 Sheet1】+ 【问题一/二输出结果】

\textbf{输出文件}: 【问题3\_图片/*.png, 问题3\_结果/*.csv】

\textbf{依赖环境}: numpy, pandas, scikit-learn, scipy

\textbf{运行顺序}: \texttt{python 问题3/问题3\_zzz.py}

\textbf{正文对应位置}: 第五章 5.3 节

% 完整可运行代码见支撑材料对应文件

\subsection*{附录 F \quad 问题四完整源程序}

\textbf{程序名称}: 问题4\_ddd.py

\textbf{编程语言及版本}: Python 3.12

\textbf{对应问题}: 问题四

\textbf{程序作用}: 【数据处理 / 模型求解 / 检验 / 绘图】

\textbf{输入文件}: 【原始数据.xlsx 的 Sheet1 或 Sheet2】+ 【问题一/二/三输出结果】

\textbf{输出文件}: 【问题4\_图片/*.png, 问题4\_结果/*.csv】

\textbf{依赖环境}: numpy, pandas, scikit-learn

\textbf{运行顺序}: \texttt{python 问题4/问题4\_ddd.py}

\textbf{正文对应位置}: 第五章 5.4 节

% 完整可运行代码见支撑材料对应文件

% 跑题后必须删下面这一行 (check 15 占位符自检)
\AppendixSlot{【这里填实际的文件清单 (附录 A 表格) + 软件环境 (附录 B) + 各问题源程序说明 (附录 C-F), 删掉这一行】}
